\begin{exercise}{10}
Show that a functor $F\colon\mathcal C\to\mathcal D$ admits a right
adjoint if you can specify objects $Gy$ for all $y\in\mathcal D$ and
maps $\varepsilon_y\colon FGy\to y$, which have the property that the
induced map on hom-sets
\[
 \operatorname{Hom}_{\mathcal C}(x,Gy)
 \xrightarrow{\ F\ }\operatorname{Hom}_{\mathcal D}(Fx,FGy)
 \xrightarrow{\ \varepsilon_y\ }\operatorname{Hom}_{\mathcal D}(Fx,y)
\]
is a bijection. Note that there is an obvious dual notion which shows
that $F$ admits a left adjoint if one can specify objects $Gy$ for all
$y\in\mathcal D$ and maps $\eta_y\colon y\to FGy$ which make the
induced map on hom-sets\footnote{The original subscripts in this
second display are $\mathcal D,\mathcal C,\mathcal C$, respectively.
They should be $\mathcal C,\mathcal D,\mathcal D$: for the given
$F\colon\mathcal C\to\mathcal D$ and $y\in\mathcal D$, the objects
$Gy,x$ lie in $\mathcal C$, while $FGy,Fx,y$ lie in $\mathcal D$.}
\[
 \operatorname{Hom}_{\mathcal C}(Gy,x)
 \xrightarrow{\ F\ }\operatorname{Hom}_{\mathcal D}(FGy,Fx)
 \xrightarrow{\ \eta_y\ }\operatorname{Hom}_{\mathcal D}(y,Fx)
\]
a bijection.

\begin{proof}
Write $\theta_{x,y}(g)=\varepsilon_yF(g)$ for the given bijections.
For $b\colon y\to y'$, define $G(b)\colon Gy\to Gy'$ as the unique
map satisfying
\begin{equation}\label{eq:adjoint-on-morphisms}
 \varepsilon_{y'}F(G(b))=b\varepsilon_y;
\end{equation}
it exists by bijectivity of $\theta_{Gy,y'}$.
The identity of $Gy$ satisfies this equation for $b=\mathrm{id}_y$,
so uniqueness gives $G(\mathrm{id}_y)=\mathrm{id}_{Gy}$.
For $c\colon y'\to y''$, we have
\[
 \varepsilon_{y''}F(G(c)G(b))
   =c\varepsilon_{y'}F(G(b))=cb\varepsilon_y,
\]
and uniqueness gives $G(cb)=G(c)G(b)$.
Thus the specified objects and these maps define a functor
$G\colon\mathcal D\to\mathcal C$.

For $a\colon x'\to x$ and $b\colon y\to y'$, equation
\eqref{eq:adjoint-on-morphisms} gives
\[
 \theta_{x',y'}(G(b)ga)=b\theta_{x,y}(g)F(a).
\]
This is naturality in both variables. The bijections $\theta_{x,y}$
therefore give the natural hom-set isomorphism required by
Definition~1.1.17, making $G$ a right adjoint of $F$.

For the dual assertion, apply the same construction to
$F^{\mathrm{op}}\colon\mathcal C^{\mathrm{op}}\to
\mathcal D^{\mathrm{op}}$, using the arrows
$\eta_y^{\mathrm{op}}\colon F^{\mathrm{op}}(Gy)\to y$.
Its required hom-set bijections are precisely those in the second
display, with arrows reversed. The resulting right adjoint of
$F^{\mathrm{op}}$ is the opposite of a left adjoint of $F$.
\end{proof}
\end{exercise}
